\documentclass[book.tex]{subfiles}

\begin{document}
\chapter{Neural Field Theory}

\label{chap:neural_quantum_fields}
\noindent\rule{11cm}{0.4pt} \\
\textbf{Prerequisites:} \autoref{chap:neural_operators}, \autoref{chap:dft},  \\
\textbf{Difficulty Level:} ** \\
\noindent\rule{11cm}{0.4pt}
\vspace{5mm}

The physical theories we've considered so far all have had finite sets of parameters. But for many more complex physical systems, we find that we're interested in modeling systems with an infinite number of parameters. Neural Field Theory provides a mechanism to directly model fields, both classical and quantum, with differentiable programs.


We define a scalar field 
\begin{align} 
\varphi: \mathbb{R}^4 \to \mathbb{R}
\end{align}
where we use $\mathbb{R}^4$ to denote spacetime. This function can be directly approximate by a suitable neural architecture
\begin{align}
\varphi &= f_\theta
\end{align}

We call such a representation a neural field.
%In \autoref{chap:maxwells_equations} we saw our first example of a field theory in this book. We will cover a number more field theories in future chapters.

\section{Electromagnetic Neural Fields}

We can represent the electric and magnetic fields $E(x)$ and $B(x)$ as neural fields. We can also represent the electromagnetic tensor $F^{\mu \nu}$ as a field. This formulation is Lorentz equivariant by construction. Maxwell's equations can be placed as a PINNs-style constraint upon the learned field.

\section{Neural Quantum Fields}

Neural quantum fields provide a technology to approximate quantum fields. Recall that a quantum field is given to first approximation by an element of type
\begin{align}
    \mathbb{R}^4_+ \to \mathcal{B}(\mathcal{H})
\end{align}
where $\mathbb{R}^4_+$ is Minkowski space. A quantum field associates to each point of Minkowski space a bounded operator on Hilbert space. 

Recall that neural operators provide a technology for approximating operators on function spaces. A neural quantum field is thus given by a function

\begin{align}
    \mathbb{R}^4_+ &\to \textrm{NeuralOperator} \\
    (x, y, z, t) &\mapsto A_{\theta, (x,y,z, t)}(\cdot)
\end{align}
where $A_{\theta, (x,y,z, t)}(\cdot)$ is a parameterized neural operator. The neural operator can be defined so that it is not dependent on the choice of discretization as we have seen in earlier chapters.

The core insight of the neural quantum field theory approach is that quantization is a much less mysterious construction in a differentiable program than in traditional physics. Quantization simply inserts a $\lambda$ (in the sense of programming language theory). This technique enables direct study of the fields with simple technology.

\subsection{Neural Distributions}
The discussion in the previous section is only an approximation, since quantum fields are more accurately described as operator valued distributions rather than operator valued functions. 

Before tackling the problem of learning an operator valued distribution, we consider the simpler problem of learning the Dirac delta. Recall that the Dirac delta is itself a distribution or a generalized function that maps an input function to its value at the origin
\begin{align}
\delta: f \mapsto f(0)
\end{align}

We can encode an approimate Dirac delta as a differentiable program
\begin{align}
\delta(f) \approx g^{\delta}_\theta 
\end{align}
where $g$ is convolution against a narrow Gaussian.

%\textcolor{red}{TODO: Explain this in more detail along with a derivation.}

\section{Klein Gordon Equation}
We can use our framework of neural operators to learn solutions to the Klein-Gordon equation. The classical Klein-Gordon equation satifies the equation

\begin{align}
(\partial_t^2 - \nabla^2 + m^2)\psi(t, x) &= 0
\end{align}

Solutions to this equation satisfy Lorentz invariance and $U(1)$ symmetry by construction. We can then quantize these solutions as in the classical solution to the quantized Klein-Gordon equation.

The usual solution to the Klein-Gordon equation converts the equation to Fourier space and notes that at each momentum value, the equation reduces to the simple harmonic oscillator. The quantization is then performed by using the ladder operators, just as the harmonic oscillator is quantized. In our case,  we can learn an approximate solution to the Klein-Gordon equation
\begin{align}
f_\theta
\end{align}
We can then use the classical commutation rules
\begin{align}
[\phi(x), \phi(y)] &= [\pi(x), \pi(y)] = 0 \\
[\phi(x), \pi(y)] &= i \delta^3(x-y)
\end{align}
to perform the commutation by the following steps. We can introduce the following forms for $\phi$ and $\pi$
\begin{align}
\phi(x) &= \lambda \psi. f_\theta(x)(\psi) \\
\pi(x) &= \lambda \psi. g_\theta(x)(\psi) \\
\end{align}
We can then introduce a loss that penalizes deviations from the commutation relations
\begin{align}
\mathcal{L}_1 &= \int_x \int_y \|[\phi(x), \phi(y)]\|^2 \\
\mathcal{L}_2 &= \int_x \int_y \|[\pi(x), \pi(y)]\|^2 \\
\mathcal{L}_3 &= \int_x \int_{y \neq x} \|[\phi(x), \phi(y)]\|^2 + \int_x \|[\phi(x), \phi(y)] - i\|^2\\
\mathcal{L} &= \mathcal{L}_1 + \mathcal{L}_2 + \mathcal{L}_3 + \mathcal{L}_4
\end{align}


%The free solution can be solved exactly (good test case?) More interesting may be to solve the solution with a field


As a more complex example, we can consider the Klein-Gordon system with an added vector potential.

\begin{align}
\Box \psi + \frac{\partial V}{\partial \psi} &= 0
\end{align}

We can use the neural field to learn solutions in cases with a vector potential just as we learned the free field above.. %The standard solution is also very hard to think about. How can that integral even be evaluated?


\section{Dirac Equation}

We can use the same methdology to construct a neural field theory for the Dirac equation

\begin{align}
i \hslash \gamma^\mu \partial_\mu \psi(x) - m c\psi(x) &= 0
\end{align}

One complication here is that $\psi(x)$ takes values in spinors, so we have to place spinorial constraint on our original PINNs-style solver. Each component of the Dirac equation satisfies the Klein-Gordon field, so we can use that constraint to help guide the solver.

We can also introduce electromagnetic coupling into the Dirac equation in a natural way, providing a pathway to apply our methodology to more complex quantum fields.

%\section{Spontaneous Symmetry Breaking}

%Relate back to the phase transition theory. How do Goldstone bosons emerge?


\section{Quantum Electrodynamics}

The QED field is a gauge field (consisting of the Dirac field plus electromagnetic coupling). The Dyson series for the perturbation expansion can be directly implemented as a differentiable program.

%\textcolor{red}{TODO: When quantizing, how do the products of operators make sense? This is a deep question I think. I think the Dyson series can be implemented as a differentiable layer.}

\section{Quantum Chromodynamics}

This same approach can be applied to the Lagrangian for quantum chromodynamics. A powerful result of this theory is that this formulation is naturally mesh-free, allowing for the natural construction of a mesh-free approximate theory of quantum chromodynamics.

%The dream here would be that neural quantum fields could allow for construction of a meshfree QCD theory. It feels like someone should have already come up with a meshfree qcd theory or at least a non-standdard mesh QCD.

%\section{Notes}

%Notes: I think the basic idea of all this is to leverage that paper that constructs Lorentz invariant quantities and construct all these fields. 

%\textcolor{red}{TODO: Can this framework handle interacting fields? To be useful in a particle accelerator, it would need to be able to handle strongly interacting fields.}

%\textcolor{red}{TODO: Should renormalizability be worked into these conditions somehow?}

%\textcolor{red}{TODO: Klein Gordon can be solved classically. How to solve as an operator? I guess the $A(p)$ and $B(p)$ are made operators.}

%\textcolor{red}{TODO: How is the field trained? Could use a PINNs like loss term?}
\end{document}